\chapter{Uncertainty Quantification}
\label{ch:uq}

\epigraph{\itshape A prediction without uncertainty is a guess. A guess
without uncertainty is malpractice.}{D.\ Saltzman}

\section{Sources of uncertainty}
\begin{itemize}
\item \textbf{Aleatoric}: irreducible randomness in the data-generating
process (turbulence, measurement noise).
\item \textbf{Epistemic}: lack of knowledge about model parameters
(reducible with more data).
\item \textbf{Structural}: wrong model class --- often the largest
contributor in SciML and the hardest to quantify.
\end{itemize}
A trustworthy SciML pipeline must acknowledge all three.

\section{Bayesian deep learning}
\subsection{Bayesian neural networks (BNNs)}
Treat weights as random variables with prior $p(\theta)$ and posterior
$p(\theta|\D)\propto p(\D|\theta)p(\theta)$. Approximations:
\begin{itemize}
\item \textbf{Variational inference}: mean-field, structured Gaussian,
normalising-flow posteriors.
\item \textbf{Laplace approximation}: $p(\theta|\D)\approx
\mathcal{N}(\theta_{\text{MAP}},\bm{H}^{-1})$ via Hessian/GGN.
\item \textbf{MCMC} (HMC, SGLD, SGNHT): asymptotically exact but costly.
\end{itemize}

\subsection{Bayesian PINNs}
B-PINN~\cite{yang2021bpinn} samples from the posterior over
$\theta$ given the residual + data likelihood, providing predictive
intervals on PDE solutions. Used for noisy inverse problems and
data-assimilation.

\section{Deep ensembles}
Train $K$ independent networks; average predictions and use the
inter-model variance as uncertainty. Often the strongest baseline ---
simple, calibrated, trivially parallel. Drawback: $K\times$ training
and inference cost. Approximations: snapshot ensembles, BatchEnsemble,
multi-input multi-output (MIMO).

\section{Gaussian processes}
The classical UQ workhorse. For a kernel $k(\bm{x},\bm{x}')$ and
observations with noise $\sigma^2$, posterior mean and variance are
analytic. Limitations: $\Oh(N^3)$ training. Modern remedies:
\begin{itemize}
\item Sparse / inducing-point methods (SVGP, FITC).
\item Stochastic variational GPs.
\item Deep kernel learning: GP head on a neural feature map.
\item Neural processes / Conditional NPs as amortised meta-GPs.
\end{itemize}
GPs are particularly natural surrogates for expensive forward simulators
and for Bayesian optimisation of physical experiments.

\section{Conformal prediction}
Distribution-free, finite-sample-valid prediction sets:
\begin{equation}
\Pr(y_{\text{new}}\in C_\alpha(\bm{x}_{\text{new}}))\ge 1-\alpha.
\end{equation}
Build a non-conformity score on a calibration set; set the prediction
interval as the $(1-\alpha)$-quantile. Variants for regression
(CQR, conformalised quantile regression), time series (split-conformal
with weighted exchangeability), and PDEs (conformalised neural
operators)~\cite{ma2024conformal}. Crucially, conformal works
\emph{on top} of any predictor and gives marginal coverage with no
distributional assumptions.

\begin{insight}{Why conformal is special}
Other UQ methods are calibrated only if the model is well-specified
(BNNs, GPs). Conformal coverage holds under exchangeability
\emph{regardless} of model correctness --- making it ideal for
hardening ML surrogates that one does not fully trust.
\end{insight}

\section{Probabilistic operators and probabilistic PINNs}
\begin{itemize}
\item \textbf{Probabilistic FNO / DeepONet}: heteroscedastic outputs;
ensemble FNOs; SDE-based latent operators.
\item \textbf{Score-based UQ}: a diffusion prior + posterior sampling
yields full predictive distributions for inverse PDE problems.
\item \textbf{Monte Carlo dropout}: cheap approximation; usually
overconfident.
\end{itemize}

\section{Calibration and evaluation}
\begin{itemize}
\item Reliability diagrams; Expected Calibration Error (ECE).
\item Continuous Ranked Probability Score (CRPS) for ensembles.
\item Negative log-likelihood (NLL) on held-out data.
\item Coverage of prediction intervals; sharpness.
\item Out-of-distribution (OOD) detection.
\end{itemize}

\section{Active learning and experimental design}
Couple UQ with sampling: select new query points to minimise expected
posterior variance (BALD, expected-information-gain). For SciML:
\begin{itemize}
\item Active learning for force-field training (NequIP-AL).
\item Bayesian optimisation for PDE-parameter calibration.
\item Multi-fidelity Bayesian optimisation when both cheap and
expensive simulators are available.
\end{itemize}

\section{Practical guide}
\begin{recipe}{UQ stack for a SciML deliverable}
\begin{enumerate}
\item Train a deterministic baseline; compute residuals on held-out data.
\item Add an ensemble of $\ge 5$ models.
\item Apply conformal prediction on a calibration split.
\item If the model will be used for inverse problems, add a Bayesian
or score-based posterior layer.
\item Report CRPS / NLL / coverage --- never just RMSE.
\end{enumerate}
\end{recipe}

\begin{pitfall}{Confusing aleatoric with epistemic}
A heteroscedastic-output network estimates aleatoric noise. Adding
ensembles or BNNs estimates epistemic. Plotting only the predictive
variance without separating the two leads to incorrect inferences ---
e.g., believing more data will reduce noise that is intrinsic to the
process.
\end{pitfall}
